\documentclass[12pt,a4paper]{article}
\usepackage{amsmath,amssymb}
\usepackage{graphicx}
\usepackage{hyperref}
\usepackage{geometry}
\usepackage{booktabs}
\usepackage{natbib}
\geometry{margin=2.5cm}
\hypersetup{colorlinks=true,linkcolor=blue,citecolor=blue,urlcolor=blue}

\newcommand{\rd}{r_{\mathrm{d}}}

\begin{document}

\title{The BAO Sound Horizon in RIFT:\\
First-Principles Derivation and Degeneracy with $\Lambda$CDM}

\author{Christopher Robert Barrick Wilson\\
\small Independent Researcher\\
\small \texttt{C.Isomorph@gmail.com}}
\date{\today}
\maketitle

\begin{abstract}
The Recursive Intelligence-Field Theory (RIFT) introduces a non-minimal
curvature coupling $\Lambda_0\Psi R$ to the scalar--tensor action.  A
longstanding open question is whether this coupling shifts the BAO comoving
sound horizon $\rd$ --- a shift that would constitute the first genuinely
distinct prediction over $\Lambda$CDM.  We derive $\rd$ directly from the
RIFT perturbation equations: the Klein--Gordon equation for $\Psi$ on the
FLRW background is integrated alongside the standard photon--baryon plasma,
and $\rd = \int_{z_{\rm drag}}^\infty c_s(z)\,c/H(z)\,dz$ is evaluated
numerically (Sim~101).  Across the full observationally allowed coupling
range $\Lambda_0\in[0,0.095]$ the maximum fractional shift is
$|\Delta\rd/\rd| = 21$~ppm, four orders of magnitude below current BAO
measurement precision ($\sim 3000$~ppm).  The scalar energy density at the
drag epoch is $\rho_\Psi/\rho_{\rm tot}\sim10^{-10}$, confirming that the
RIFT background is indistinguishable from FLRW at $\Lambda_0\ll 1$.
The echo-shell mechanism proposed in the parent RIFT paper is therefore
a physical reinterpretation of standard acoustic physics, not a derivation
of a distinct $\rd$.  The detection threshold for a RIFT sound-horizon
signal requires $\Lambda_0\gg 0.1$, well beyond the current 95\% upper limit.
\end{abstract}

\section{Introduction}

Baryon Acoustic Oscillations fix a standard ruler at the comoving sound
horizon
\begin{equation}
\rd = \int_{z_{\rm drag}}^{\infty} \frac{c_s(z)\,c}{H(z)}\,dz\,,
\label{eq:rd}
\end{equation}
where $c_s = c/\sqrt{3(1+R_b)}$ is the photon--baryon sound speed,
$R_b = 3\rho_b/(4\rho_\gamma)$, and the drag redshift $z_{\rm drag}\approx1059$
marks the decoupling of baryons from radiation.  In $\Lambda$CDM the
integral is dominated by the radiation era and yields
$\rd\approx 147$~Mpc (Planck~2018).

RIFT extends the Einstein--Hilbert action with a non-minimal coupling
$\Lambda_0\Psi R$, where $\Psi$ satisfies the Klein--Gordon equation
\begin{equation}
\ddot\Psi + 3H\dot\Psi + m_0^2\Psi = -2\Lambda_0 R\,.
\label{eq:KG}
\end{equation}
If $\Psi$ carries non-negligible energy density during recombination,
or if the RIFT modification of $H(z)$ is significant, $\rd$ would shift.
This paper quantifies that shift from first principles.

\section{Method: RIFT Perturbation Equations at Recombination}

We integrate the background FLRW system numerically from $z=10^6$ to $z=0$:
\begin{align}
H^2(z) &= H_0^2\bigl[\Omega_m(1+z)^3 + \Omega_r(1+z)^4
  + \Omega_\Lambda + \Omega_\Psi(z)\bigr]\,,\\
\frac{du}{dz} &= u\!\left(\frac{2}{1+z} - \frac{H'}{H}\right)
  - \frac{m_0^2\Psi}{(1+z)^2 H^2}\,,\quad u \equiv \frac{d\Psi}{dz}\,,
\end{align}
where $\Omega_\Psi(z) = (\dot\Psi^2/2 + m_0^2\Psi^2/2)/(3H_0^2M_{\rm Pl}^2)$.
The sound horizon is computed via eq.~\eqref{eq:rd} with CLASS-calibrated
$z_{\rm drag}=1059$.

Best-fit parameters from the joint CMB+BAO fit (Sim~90):
$H_0=67.59$~km/s/Mpc, $\Omega_m=0.312$, $\Omega_r=9.24\times10^{-5}$,
$m_0=0.01$~Mpc$^{-1}$, $\Psi_0(z=0)=0.003$.

\section{Results}

\begin{table}[h]
\centering
\caption{BAO sound horizon versus RIFT coupling.}
\label{tab:rd}
\begin{tabular}{cccc}
\toprule
$\Lambda_0$ & $\rd^{\rm RIFT}$ [Mpc] & $\Delta\rd$ [Mpc] & $\Delta\rd/\rd$ \\
\midrule
$0$ ($\Lambda$CDM) & $153.05$ & $0$ & $0$ \\
$0.003$ & $153.05$ & $-1.0\times10^{-4}$ & $-0.68$~ppm \\
$0.008$ & $153.05$ & $-2.8\times10^{-4}$ & $-1.8$~ppm \\
$0.050$ & $153.05$ & $-1.7\times10^{-3}$ & $-11$~ppm \\
$0.095$ & $153.05$ & $-3.3\times10^{-3}$ & $-21$~ppm \\
\bottomrule
\end{tabular}
\end{table}

The scalar energy density at the drag epoch is
$\rho_\Psi/\rho_{\rm tot}\approx 1.5\times10^{-10}$ for all scanned
$\Lambda_0$, confirming the RIFT background is indistinguishable from
FLRW at recombination.

The $\Lambda$CDM baseline $\rd = 153$~Mpc differs from the Planck value
of 147~Mpc because the simplified Friedmann equation omits neutrino
free-streaming ($\Delta N_{\rm eff}$); both RIFT and $\Lambda$CDM share
this approximation, so the fractional shifts are unaffected.

\section{Discussion}

The maximum shift $|\Delta\rd/\rd| = 21$~ppm is four orders of magnitude
below the $\sim 3000$~ppm (0.3\%) measurement precision of current BAO
surveys.  RIFT and $\Lambda$CDM are degenerate in $\rd$ at all
observationally allowed couplings.

The echo-shell picture advanced in the parent RIFT paper --- that the
$\Lambda_0\Psi R$ coupling reinforces acoustic oscillations and shifts
the BAO peak --- is confirmed as a physical interpretation of standard
photon--baryon acoustics rather than a first-principles derivation
of a distinct scale.  The sound horizon $\rd$ remains a free parameter
fitted to data; RIFT does not predict a different value from $\Lambda$CDM
within the locked action.

The detection threshold for a RIFT $\rd$ signal requires
$\Lambda_0\gg 0.1$, well beyond the current 95\% upper limit
$\Lambda_0<0.095$ from joint CMB+BAO fitting (Sim~90, 93).
Next-generation surveys (DESI~Y5, Euclid) will constrain $\rd$ to
$\sim 0.1\%$, probing $\Lambda_0\gtrsim 0.1$ --- just at the edge of
the current bound.

\section{Conclusion}

A direct first-principles derivation of the BAO sound horizon in RIFT
yields $|\Delta\rd/\rd|<21$~ppm across all allowed couplings.
The locked action at $\Lambda_0\approx 0.003$ does not predict a
measurably different $\rd$ from $\Lambda$CDM.
This closes the P2.7 referee concern raised for the parent paper:
RIFT is consistent with all BAO data not because it predicts the
same ruler, but because it \emph{is} the same ruler at current couplings.
A detectable $\rd$ shift is a falsifiable prediction for
$\Lambda_0\gtrsim 0.1$.

\end{document}
